% #############################################################################
% This is Chapter 6
% !TEX root = ../main.tex
% #############################################################################
% Change the Name of the Chapter i the following line
\fancychapter{Conclusion}
% The following line allows to ref this chapter
\label{chap:conclusion}

This chapter concludes the thesis. On the first section, a summary of the main points discussed in the previous chapters is presented, along with the contributions of this work. The second section outlines potential areas for future work and improvements. 

\section{Conclusions}

In this thesis, a novel solution for compiling a linear session-typed language, \ac{CLASS} \cite{Rocha:2023}, into efficient C code was designed and implemented. The first \acf{IR} for a session-typed language, suitable for optimizations and easily translatable from session-typed languages and to efficient low-level imperative code, was also developed. The solution preserves the semantics and safety guarantees of the source language, while generating code that is orders of magnitude faster than the interpreted execution of \ac{CLASS} programs. The compilation scheme is based on the sequential execution model of \ac{CLASS} used by the \acf{SAM} \cite{Caires:2024}.

The developed solution consists of a type-directed compilation algorithm which translates \ac{CLASS} programs into \ac{IR} programs, on which several optimizations are performed, before finally being translated into C code. The compilation makes use of the information available in the session types of \ac{CLASS} to generate efficient code, avoiding unnecessary runtime checks and garbage collection.

All program data, excluding exponential closures and shared state cells, is stored in structures associated to processes, similar to the stack frames used in imperative languages, but where the frames are organized into a spaghetti stack. Instead of the traditional single return address, each frame stores a continuation for each unfinished session in the process.

Communication over sessions is compiled into direct memory writes reads into and from these frames, with sizes and offsets determined at compile time, excluding polymorphic code sections. When a sequence of writes on a given session finishes, control is passed back to the continuation stored in the active frame, following the execution model of the \ac{SAM}.

An implementation of the developed scheme was created, covering all of the features of \ac{CLASS}, including polymorphism, recursion and exponentials. The implementation, available on GitHub \footnote{\url{https://github.com/RiscadoA/class}}, was built in Java, on top of the existing \ac{SAM} codebase. A suite of tests was created to validate the correctness of the implementation, and the performance of the generated code was evaluated through a series of benchmarks, across different configurations. The compiler was shown to generate programs competitive with equivalent Haskell GHC compiled programs, being able to surpass them in some cases.

% #############################################################################
\section{Future Work}

The developed solution, while already showing promising results, still has several avenues for future improvements and optimizations. One such avenue is compiling to a lower level target language, such as LLVM IR, or assembly language. This could allow for, for example, storing frequently accessed data entirely in CPU registers, and committing them to memory only when absolutely necessary.

The \ac{IR} could also be optimized further, by removing, for example, redundant moves and copies of data, which are frequently generated by the current compilation scheme. Both sessions and local data sections could be reused if possible, reducing the memory footprint of the generated programs, and improving cache locality.

It would also be interesting to construct a formal proof of correctness for the compilation scheme, demonstrating that the generated \ac{IR} code preserves the semantics and safety of \ac{CLASS}. As the compiled code strongly resembles the \ac{SAM} execution model, the proof could be based on its existing proof of correctness.

Concurrent sessions have not been addressed in this work, as the focus was on generating efficient sequential code. The type system of the \ac{CLASS} language could be modified to distinguish between sequential and concurrent sessions, allowing for the generation of efficient code for both sessions types.

Finally, a standard library for \ac{CLASS} could also be developed, providing native implementations of common data structures such as arrays and maps, as currently these are emulated using recursive session types. This would greatly improve the usability and performance of \ac{CLASS} programs. This could be achieved by, for example, implementing a \ac{FFI} to C libraries.
